% Lab 5 -- Exercise 1  |  Nguyễn Huỳnh Hữu Đức – 2510290  |  Original: Ch.6 2.1 (JFET DC sweep)  |  Status: Adjust
% Netlist để vẽ schematic (LabN/Ex01.asc):
%   V2  d 0 10          ; nguồn drain
%   V3  g 0 0           ; nguồn gate (được quét)
%   J1  d g 0 JN        ; JFET kênh N, source nối đất
\exercise{JFET Transfer Characteristic}{Adjust}{Ch.6 2.1 -- JFET DC sweep}
\statusreason{In the original manual this sweep is a demonstration with the PSpice part \texttt{JbreakN}.
It is turned into an exercise with a theoretical part, and the device is defined with an explicit LTspice
\texttt{.model} ($V_P=-2$\,V, $I_{DSS}=4$\,mA) so that students can predict the curve with Shockley's equation.
The learning objective (reading $I_{DSS}$ and $V_P$ from a DC sweep) is unchanged.}

\task{An n-channel JFET is biased with a fixed drain--source voltage $V_{DS} = 10$\,V. The gate--source voltage
$V_{GS}$ is swept from $-3$\,V to $0$\,V. Predict the drain current with Shockley's equation, then plot the
transfer characteristic $I_D(V_{GS})$ in LTspice and read $I_{DSS}$ and $V_P$ from the curve.}

\subsubsection*{Circuit}
\circuitfig{ex01}{Circuit of Exercise 5.1: JFET transfer characteristic}
\components{
\begin{itemize}[leftmargin=*,itemsep=1pt]
  \item $V_2 = 10$\,V between drain and ground; $V_3$ (swept) between gate and ground; source grounded.
  \item J1: n-channel JFET (symbol \texttt{njf}), model \texttt{JN}:\\
        \texttt{.model JN NJF(Vto=-2 Beta=1m Lambda=0)}
\end{itemize}}

\subsubsection*{Theoretical Calculation}
\begin{enumerate}[itemsep=2pt]
  \item From the model, find $V_P$ and $I_{DSS}$. (In LTspice, $V_P$ = \texttt{Vto} and $I_{DSS} = \texttt{Beta}\cdot V_P^2$.)
  \item Use Shockley's equation $I_D = I_{DSS}\left(1 - V_{GS}/V_P\right)^2$ to compute $I_D$ at
        $V_{GS} = 0$, $-0.5$, $-1$, $-1.5$ and $-2$\,V.
\end{enumerate}
\answerspace{8}
\begin{solution}
$V_P = -2$\,V and $I_{DSS} = \beta V_P^2 = (1\,\mathrm{mA/V^2})(4\,\mathrm{V^2}) = 4$\,mA.
\[
  I_D = 4\,\mathrm{mA}\left(1+\frac{V_{GS}}{2}\right)^2
\]
\begin{center}
\begin{tabular}{c c c c c c}
\toprule
$V_{GS}$ (V) & 0 & $-0.5$ & $-1$ & $-1.5$ & $-2$\\ \midrule
$I_D$ (mA) & 4.00 & 2.25 & 1.00 & 0.25 & 0\\
\bottomrule
\end{tabular}
\end{center}
\end{solution}

\subsubsection*{Simulation}
\begin{enumerate}[itemsep=2pt]
  \item Draw the circuit, place the \texttt{.model} line as a SPICE directive and set the model name of J1 to \texttt{JN}.
  \item Add the DC sweep below and run. Plot \texttt{Id(J1)}.
  \item Use the cursor to read $I_D$ at $V_{GS}=0$ ($I_{DSS}$), the voltage where $I_D$ becomes zero ($V_P$), and $I_D$ at $V_{GS}=-1$\,V.
\end{enumerate}
\begin{spice}
.model JN NJF(Vto=-2 Beta=1m Lambda=0)
.dc V3 -3 0 0.01
\end{spice}
\simfig{ex01_sim}{Transfer characteristic $I_D(V_{GS})$ of Exercise 5.1}

\subsubsection*{Theory vs. Simulation}
\begin{comparison}{Theory vs. simulation, Exercise 5.1}
\cmprow{$I_{DSS}$}{mA}{4.00}{}{}
\cmprow{$V_P$}{V}{$-2.00$}{}{}
\cmprow{$I_D$ at $V_{GS}=-1$\,V}{mA}{1.00}{}{}
\cmprow{$I_D$ at $V_{GS}=-1.5$\,V}{mA}{0.25}{}{}
\end{comparison}
\answerspace{3}
\begin{solution}
With \texttt{Lambda=0} the model follows Shockley's equation exactly, so the deviation should be limited to the cursor
resolution (step 10\,mV). \todo{Fill in the simulation column from the LTspice run.}
\end{solution}
